\documentclass[10pt,a4paper]{amsart}
\usepackage{amsmath,amssymb,amsthm,graphicx}
% Retained minimal preamble for the Stacks Project chapter more-algebra.tex (More on Algebra).
% The chapter file pinned-original-source.tex carries only \input{preamble} on its first line, so no
% preamble text lives inside the chapter itself.  The Stacks Project uses one single shared preamble for
% every chapter; the snapshot stages it identically as sources/stacks-chow/preamble.tex and
% sources/stacks-etale-cohomology/preamble.tex.  The block below is copied mechanically, line for line,
% from sources/stacks-chow/preamble.tex lines 14-18 (the Xy-pic setup) and 154-217 (the native theorem
% declarations together with the mathematical macros).  Nothing was added, removed or redefined; the
% wrapper class and the amsmath/amssymb/amsthm/graphicx base packages are supplied by the assembler.
% Chapter: more-algebra.tex -- More on Algebra.
% For commutative diagrams we use Xy-pic
\usepackage[all]{xy}

% We use 2cell for 2-commutative diagrams.
\xyoption{2cell}

% Theorem environments.
%
\theoremstyle{plain}
\newtheorem{theorem}[subsection]{Theorem}
\newtheorem{proposition}[subsection]{Proposition}
\newtheorem{lemma}[subsection]{Lemma}

\theoremstyle{definition}
\newtheorem{definition}[subsection]{Definition}
\newtheorem{example}[subsection]{Example}
\newtheorem{exercise}[subsection]{Exercise}
\newtheorem{situation}[subsection]{Situation}

\theoremstyle{remark}
\newtheorem{remark}[subsection]{Remark}
\newtheorem{remarks}[subsection]{Remarks}

\numberwithin{equation}{subsection}

% Macros
%
\def\lim{\mathop{\mathrm{lim}}\nolimits}
\def\colim{\mathop{\mathrm{colim}}\nolimits}
\def\Spec{\mathop{\mathrm{Spec}}}
\def\Hom{\mathop{\mathrm{Hom}}\nolimits}
\def\Ext{\mathop{\mathrm{Ext}}\nolimits}
\def\SheafHom{\mathop{\mathcal{H}\!\mathit{om}}\nolimits}
\def\SheafExt{\mathop{\mathcal{E}\!\mathit{xt}}\nolimits}
\def\Sch{\mathit{Sch}}
\def\Mor{\mathop{\mathrm{Mor}}\nolimits}
\def\Ob{\mathop{\mathrm{Ob}}\nolimits}
\def\Sh{\mathop{\mathit{Sh}}\nolimits}
\def\NL{\mathop{N\!L}\nolimits}
\def\CH{\mathop{\mathrm{CH}}\nolimits}
\def\proetale{{pro\text{-}\acute{e}tale}}
\def\etale{{\acute{e}tale}}
\def\QCoh{\mathit{QCoh}}
\def\Ker{\mathop{\mathrm{Ker}}}
\def\Im{\mathop{\mathrm{Im}}}
\def\Coker{\mathop{\mathrm{Coker}}}
\def\Coim{\mathop{\mathrm{Coim}}}

% Boxtimes
%
\DeclareMathSymbol{\boxtimes}{\mathbin}{AMSa}{"02}

%
% Macros for moduli stacks/spaces
%
\def\QCohstack{\mathcal{QC}\!\mathit{oh}}
\def\Cohstack{\mathcal{C}\!\mathit{oh}}
\def\Spacesstack{\mathcal{S}\!\mathit{paces}}
\def\Quotfunctor{\mathrm{Quot}}
\def\Hilbfunctor{\mathrm{Hilb}}
\def\Curvesstack{\mathcal{C}\!\mathit{urves}}
\def\Polarizedstack{\mathcal{P}\!\mathit{olarized}}
\def\Complexesstack{\mathcal{C}\!\mathit{omplexes}}
% \Pic is the operator that assigns to X its picard group, usage \Pic(X)
% \Picardstack_{X/B} denotes the Picard stack of X over B
% \Picardfunctor_{X/B} denotes the Picard functor of X over B
\def\Pic{\mathop{\mathrm{Pic}}\nolimits}
\def\Picardstack{\mathcal{P}\!\mathit{ic}}
\def\Picardfunctor{\mathrm{Pic}}
\def\Deformationcategory{\mathcal{D}\!\mathit{ef}}

% --- appended 2026-10-12 (Chief): the project own definitions of the reference/slogan/history
% display environments, copied mechanically from the shared Stacks preamble lines 8-12
% (\usepackage{verbatim} plus the three \newenvironment lines, which the chapter preamble keeps
% verbatim from the project preamble). Nothing added, removed or redefined: these environments
% carry no mathematical content and the project itself typesets them as comments. Without these
% three lines a counted statement whose own text uses them fails with Environment slogan undefined.
\usepackage{verbatim}
\newenvironment{reference}{\comment}{\endcomment}
\newenvironment{slogan}{\comment}{\endcomment}
\newenvironment{history}{\comment}{\endcomment}
\usepackage{cleveref}
\crefname{originallemma}{original lemma}{original lemmas}
\Crefname{originallemma}{Original lemma}{Original lemmas}
\makeatletter
\newlabel{more-algebra-lemma-extended-alternating-torsion}{{0G6K}{source}{}{sourcepointer-1}{}}
\newlabel{more-algebra-lemma-extended-alternating-torsion@cref}{{[originallemma][0][]0G6K}{[1][0][]source}{}{}{}}
\makeatother
\crefname{originallemma}{original lemma}{original lemmas}
\Crefname{originallemma}{Original lemma}{Original lemmas}
\makeatletter
\newlabel{lemma-extended-alternating-base-change}{{0G6I}{source}{}{sourcepointer-2}{}}
\newlabel{lemma-extended-alternating-base-change@cref}{{[originallemma][0][]0G6I}{[1][0][]source}{}{}{}}
\makeatother
\crefname{originallemma}{original lemma}{original lemmas}
\Crefname{originallemma}{Original lemma}{Original lemmas}
\makeatletter
\newlabel{lemma-extended-alternating-homotopy-zero}{{0G6J}{source}{}{sourcepointer-3}{}}
\newlabel{lemma-extended-alternating-homotopy-zero@cref}{{[originallemma][0][]0G6J}{[1][0][]source}{}{}{}}
\makeatother
\title{Torsion and vanishing properties of the cohomology of the extended alternating Cech complex}
\author{The Stacks Project Authors}
\date{stacks-project:more-algebra.tex:local-snapshot-20261012; GNU FDL 1.2 or later}
\begin{document}
\maketitle
\setcounter{subsection}{1}
\noindent\textit{Editorial note.} All delivered mathematical text is the author's own, copied from the pinned Stacks Project chapter file more-algebra.tex; no statement or proof step was rewritten, reordered or removed.\par\medskip
\noindent\textit{Editorial note.} The chapter file pinned-original-source.tex carries only \textbackslash{}input\{preamble\} on its first line, so no preamble text lives inside the chapter. The retained file sources/stacks-more-algebra/retained-preamble.tex is a mechanical, line-for-line copy of the project's single shared preamble as staged in the snapshot (Xy-pic setup plus theorem declarations and mathematical macros). Nothing was added, removed or redefined, and the excerpt serves only as preamble, never as a span coordinate.\par\medskip
\noindent\textit{Editorial note.} Both counted spans start at the \textbackslash{}begin line of their environment, so the statement and the author's proof are delivered whole; surrounding prose of the section is deliberately outside the counted spans and is carried as standard prerequisites instead.\par\medskip
\noindent\textit{Editorial note.} The citation fields (source\_id, source\_version, authors, source\_url, license.name, source\_title) are identical to those of every other card this lane has produced from this chapter, so the batch selector treats them as one source.\par\medskip
\noindent\textit{Allowed standard prerequisites (editorial).}\begin{itemize}
\item Extended alternating Cech complex of a sequence and its shape in each degree. recalled in the introduction of the same block, source lines 6908-7064. Explains why the complex vanishes outside the range of degrees.
\item Base change of the extended alternating Cech complex: localising the complex at an element is the complex of the localised module over the localised ring. the statement counted by the sibling card S1-algebra-extended-alternating-Cech-is-colimit-koszul at 7099-7197 and the surrounding block; the author cites Lemma lemma-extended-alternating-base-change at proof lines 7086-7090. Identifies the localised cohomology with the cohomology of the localised complex.
\item The extended alternating Cech complex is homotopy equivalent to zero when one of the elements is a unit. the statement counted by the sibling card S1-algebra-extended-alternating-homotopy-zero at 6980-7063, cited at proof lines 7090-7092. Supplies the vanishing that the annihilation statement reduces to.
\end{itemize}\medskip
\ifdefined\hypertarget\hypertarget{sourcepointer-1}{}\fi
\noindent\textit{Source reference.} Original lemma 0G6K: focused-seeds/AGP/sources/stacks-snapshot-20261011/tags.txt, source lines 20200--20200. The lemma's own canonical tag, quoted so that readers can resolve the label inside the delivered text; it points at the counted statement itself and carries no proof content.\par\medskip
\ifdefined\hypertarget\hypertarget{sourcepointer-2}{}\fi
\noindent\textit{Source reference.} Original lemma 0G6I: focused-seeds/AGP/sources/stacks-snapshot-20261011/tags.txt, source lines 20198--20198. Navigation pointer to a statement the counted proof refers to but does not contain; nothing load-bearing is delegated to it.\par\medskip
\ifdefined\hypertarget\hypertarget{sourcepointer-3}{}\fi
\noindent\textit{Source reference.} Original lemma 0G6J: focused-seeds/AGP/sources/stacks-snapshot-20261011/tags.txt, source lines 20199--20199. Navigation pointer to a statement the counted proof refers to but does not contain; nothing load-bearing is delegated to it.\par\medskip
\begin{lemma}
\label{lemma-extended-alternating-torsion}
Let $R$ be a ring. Let $f_1, \ldots, f_r \in R$. Let $M$ be an $R$-module.
Let $H^q$ be the $q$th cohomology module of the extended alternation
{\v C}ech complex of $M$. Then
\begin{enumerate}
\item $H^q = 0$ if $q \not \in [0, r]$,
\item for $x \in H^i$ there exists an $n \geq 1$ such that $f_i^n x = 0$
for $i = 1, \ldots, r$,
\item the support of $H^q$ is contained in $V(f_1, \ldots, f_r)$,
\item if there is an $f \in (f_1, \ldots, f_r)$ which acts invertibly
on $M$, then $H^q = 0$.
\end{enumerate}
\end{lemma}

\begin{proof}
Part (1) follows from the fact that the extended alternating {\v C}ech complex
is zero in degrees $< 0$ and $> r$. To prove (2) it suffices to show that
for each $i$ there exists an $n \geq 1$ such that $f_i^n x = 0$. To see
this it suffices to show that $(H^q)_{f_i} = 0$. Since localization is
exact, $(H^q)_{f_i}$ is the $q$th cohomology module of the localization
of the extended alternating complex of $M$ at $f_i$. By
Lemma \ref{lemma-extended-alternating-base-change}
this localization is the extended alternating {\v C}ech complex
of $M_{f_i}$ over $R_{f_i}$ with respect to the images of
$f_1, \ldots, f_r$ in $R_{f_i}$. Thus
we reduce to showing that $H^q$ is zero if $f_i$ is invertible, which
follows from Lemma \ref{lemma-extended-alternating-homotopy-zero}.
Part (3) follows from the observation that $(H^q)_{f_i} = 0$
for all $i$ that we just proved. To see part (4) note that in this
case $f$ acts invertibly on $H^q$ and $H^q$ is supported on $V(f)$
by (3). This forces $H^q$ to be zero (small detail omitted).
\end{proof}

\end{document}
